\section{Conclusion}
\label{sec:conclusion}

We showed that the standard post-training recipe of a world action model breaks down when only three demonstrations per task are available, losing 16 to 49 points on the four LIBERO suites, and traced the problem to a mismatch between the objective of mid-training, which seeks weights that are good for a mixture of embodiments, and the need of few-shot post-training, which requires weights from which a short fine-tune on a new embodiment succeeds.
\Method closes part of this gap with a Reptile stage that repeatedly simulates few-shot post-training on source embodiments already present in mid-training and moves the entire trainable set of the recipe toward the adapted solutions.
The stage needs no new data, runs on the unchanged sharded training stack, and produces a checkpoint that the unchanged post-training recipe consumes.
It improves all four suites in the few-shot regime, every task of LIBERO-10, and the full-data setting as well.
Separating the meta-learning rule from the extra data exposure, meta-training with multi-instruction episodes to address instruction grounding, and scaling the stage to larger backbones are the next steps.
